\section{Flatten a Binary Tree to a Linked List}
\label{sec:trees-flatten}

\problemstatement{Flatten a binary tree in place so that it becomes a
``linked list'': every node's left child is null and its right child points
to the next node in \emph{preorder}. Use the existing right pointers as the
list's next pointers.}

\begin{patternbox}
The target order is \emph{preorder}, so the flattened list is node, then
all of the (flattened) left subtree, then all of the (flattened) right
subtree. The elegant $O(1)$-space solution is a \textbf{reverse-preorder}
(right, left, node) walk that stitches each node's right pointer to the
previously visited node -- building the list backwards.
\end{patternbox}

\subsection*{Understanding the problem carefully}

In the final list, after any node come all its left-subtree nodes, then its
right-subtree nodes. If we process nodes in the \emph{reverse} of preorder
-- that is right subtree, then left subtree, then the node -- then at the
moment we handle a node, the entire portion of the list that should come
\emph{after} it has already been assembled and is pointed to by a running
\textit{prev} pointer. So we set node.right $=$ \textit{prev}, node.left
$=$ null, and update \textit{prev} $=$ node.

\begin{ideabox}[Build the List Backwards with a Running \textit{prev}]
Traverse in right, left, node order. Keep \textit{prev}, the head of the
already-flattened suffix. At each node: set its right pointer to
\textit{prev}, its left pointer to null, then set \textit{prev} to this
node. When the traversal finishes, \textit{prev} is the root of the fully
flattened list.
\end{ideabox}

\subsection*{Algorithm}

\begin{enumerate}
  \item \textit{prev} $=$ null.
  \item \textsc{Flatten}(node):
  \begin{enumerate}
    \item If node is null, return.
    \item \textsc{Flatten}(node.right).
    \item \textsc{Flatten}(node.left).
    \item node.right $=$ \textit{prev}; node.left $=$ null; \textit{prev}
          $=$ node.
  \end{enumerate}
\end{enumerate}

\subsection*{Dry run}

\dryrun{Tree: root $1$; left $2$ (children $3,4$); right $5$ (right child
$6$). Preorder: $1,2,3,4,5,6$.}

\begin{itemize}
  \item Reverse-preorder processes $6$ first: $6$.right$=$null,
        \textit{prev}$=6$.
  \item Then $5$: $5$.right$=6$, \textit{prev}$=5$. Then $4$:
        $4$.right$=5$, \textit{prev}$=4$. Then $3$: $3$.right$=4$. Then
        $2$: $2$.right$=3$, left null. Then $1$: $1$.right$=2$, left null.
  \item Result chain via right pointers: $1\!\to\!2\!\to\!3\!\to\!4\!\to
        \!5\!\to\!6$, all left pointers null.
\end{itemize}

\subsection*{C++ implementation}

\begin{lstlisting}
#include <bits/stdc++.h>
using namespace std;

struct TreeNode {
    int val;
    TreeNode* left;
    TreeNode* right;
    TreeNode(int v) : val(v), left(nullptr), right(nullptr) {}
};

TreeNode* prevNode = nullptr;

void flatten(TreeNode* node) {
    if (node == nullptr) return;
    flatten(node->right);      // right first
    flatten(node->left);       // then left
    node->right = prevNode;    // stitch onto the already-built suffix
    node->left = nullptr;
    prevNode = node;
}

int main() {
    TreeNode* root = new TreeNode(1);
    root->left = new TreeNode(2);
    root->right = new TreeNode(5);
    root->left->left = new TreeNode(3);
    root->left->right = new TreeNode(4);
    root->right->right = new TreeNode(6);

    prevNode = nullptr;
    flatten(root);

    for (TreeNode* c = root; c != nullptr; c = c->right)
        cout << c->val << " ";
    cout << "\n";
    return 0;
}
\end{lstlisting}

\begin{complexitybox}
\textbf{Time Complexity.} $O(n)$ -- each node is visited once.
\textbf{Space Complexity.} $O(h)$ for the recursion stack; a Morris-style
variant achieves $O(1)$ by rewiring the rightmost node of each left subtree
to the right subtree, in the spirit of the previous two sections.
\end{complexitybox}

\begin{takeawaybox}
Processing nodes in reverse of the target order, while carrying a
\textit{prev} pointer to the already-built tail, is a recurring in-place
linking technique -- the same idea that reverses a linked list. Flattening
a tree is that trick applied along a preorder spine.
\end{takeawaybox}
